\scanpage{5}\mspage{4}

Picard - on peut dire que la classification ``au sens Picard'' est plus fine que la classification au sens des (Gr)-catégories. \deleted{\unsure{x}}

Associant à tout $X\handwritten{\in}\mathrm{Ob}\ C$ la symétrie de $X\otimes X$, interprété comme \deleted{\rlap{xxxx}étan} étant la tensorisation par un élément bien déterminé $s(X)$ de $\handwritten{\pi}_{1}(C)=\mathrm{Aut}(1_C)$, on trouve un\deleted{e} \deleted{\rlap{xxxxxxxxxxx}application}\inserted{homomorphisme} canonique
\[
s(C)=s:\handwritten{\pi}_{\deleted{1}\handwritten{0}}\handwritten{\longrightarrow}\unsure{{}'|\cdot}{}_{2}(\handwritten{\pi}_{\deleted{0}\handwritten{1}})\ ,
\]
où $ {}_{2}(\handwritten{\pi}_{\deleted{0}1})$ désigne le sous-groupe des éléments $x$ de $\handwritten{\pi}_{1}$ tels que $2x=0$.

Le résultat principal de la théorie de structure des catégories de Picard peut s'énoncer en disant que $C$ est déterminée (à équivalence \inserted{\handwritten{non canonique}} de catégories de Picard près) par la donnée des groupes commutatifs $\handwritten{\pi}_{0}$ et $\handwritten{\pi}_{1}$, et l'homomorphisme $s:\handwritten{\pi}_{0}\handwritten{\longrightarrow}{}_{2}(\handwritten{\pi}_{1})$. Dire que cet homomorphisme est nul signifie d'ailleurs que $C$ est une catégorie de Picard \uline{stricte}, et on retrouve le résultat de \handwritten{[\, ]} suivant lequel une catégorie de Picard stricte est déterminée (à équivalence près) par la connaissance de ses seuls invariants $\handwritten{\pi}_{0}$ et $\handwritten{\pi}_{1}$.

\uline{Remarques}. La théorie de loc. cit. nous permet, plus précisément, d'interpréter la \deleted{e}2-catégorie des catégories de \deleted{E}Picard strictes en termes de la sous-catégorie pleine de la catégorie dérivée de (Ab) (cat. des groupes abéliens), formée des complexes $K^{\deleted{\unsure{\bullet|\cdot}}}_{\handwritten{\bullet}}$ tels que $H^{\deleted{\unsure{i|1}}}_{\handwritten{i}}(K^{\deleted{\bullet}}_{\handwritten{\bullet}})=0$ pour $i\handwritten{\ne}0,1$. Le résultat précédent, assez grossier, se réduit \inserted{alors} à la remarque qu'un objet de cette catégorie est déterminé \deleted{\unsure{à|a}}à isomorphisme (non unique \handwritten{!}) près par la connaissance de son $H^{\deleted{1}}_{\handwritten{0}}$ $(=\handwritten{\pi}_{0})$ et $H^{\deleted{0}}_{\handwritten{1}}$ $(=\handwritten{\pi}_{1})$. Il y aurait lieu de développer également pour les catégories de Picard non nécessairement strictes, et pour les (Gr)-catégories, une théorie de structure pour la 2-catégorie qu'ils forment. De plus, ces \deleted{\rlap{xxxxxxxxx}résultats}\inserted{développements} encore dans les limbes, \deleted{\unsure{de|ce|co}}tout comme ceux que \deleted{\rlap{xxxxxxxxxxxxxxxxxxxxxxxx}nous avons effectivement} nous venons d'esquisser, devraient être un jour \deleted{\rlap{xxxxxxxxxx}étendus \unsure{éu|áu|au}} étudiés dans le cadre des champs (en (Gr)-catégories, ou en catégories de Picard) tout comme dans \inserted{\handwritten{Deligne}} \handwritten{[\, ]}.

\deleted{2} \handwritten{3}. \uline{Catégories de Picard enveloppantes}.

Soient $C$ une \deleted{\rlap{xx}\unsure{ca|la}} $\otimes$-catégorie AUC (par ëxemple une catégorie avec objet final et admettant des \inserted{\handwritten{\unsure{''|:}}}produits cartésiens de deux éléments), $S$ une partie de $\mathrm{Ob}\ C$. On se pose le problème universel d'envoyer (par un $\otimes$-foncteur AUC) $C$ dans une \inserted{$\otimes$-}catégorie \deleted{\unsure{$G$}}\inserted{AUC $G$}, de tel\handwritten{\unsure{é|le}} façon que pour tout $X\handwritten{\in} S$, $F(X)$ soit un objet \uline{inversible} de $G$. Le résultat principal ici est que ce problème admet toujours une solution, qu'on pourra appeler la $\otimes$-catégorie AUC ``de fractions'' par rapport à l'ensemble $S$, et noter \handwritten{\unsure{é|e}}ventuellement $S^{-1}C$. On trouve en fait que pour toute $G$, le foncteur naturel
% archive sheet number: 4
